\section{RAG và AI Agent}

\subsection{Retrieval-Augmented Generation}

Retrieval-Augmented Generation (RAG) kết hợp truy xuất tài liệu với mô hình
ngôn ngữ. Câu hỏi được dùng để tìm các đoạn liên quan; nội dung truy xuất được
đưa vào ngữ cảnh để tạo câu trả lời có căn cứ \cite{lewis2020rag}. Trong đồ án,
Knowledge chứa runbook, hướng dẫn triển khai, mô tả lỗi và tài liệu hệ thống.

\subsection{Embedding Search}

Embedding search ánh xạ truy vấn và tài liệu thành vector rồi xếp hạng theo độ
tương đồng. Phương pháp này tìm được nội dung gần nghĩa dù từ khóa không trùng
khớp. Chất lượng phụ thuộc cách chia đoạn, mô hình embedding, metadata filter và
ngưỡng lấy kết quả.

\subsection{Hybrid Search}

Hybrid search kết hợp tìm kiếm từ khóa với vector search. Từ khóa phù hợp với mã
lỗi, model version và tên node; vector search phù hợp với câu hỏi diễn đạt tự
nhiên. Hai danh sách có thể được chuẩn hóa điểm, hợp nhất và rerank trước khi đưa
cho Agent.

\subsection{AI Agent}

AI Agent mở rộng RAG bằng khả năng lập kế hoạch và gọi công cụ. Agent có thể
truy vấn Fleet, Runtime, Deployment hoặc Incident thay vì chỉ trả lời từ tài liệu.
Mô hình ngôn ngữ không được truy cập trực tiếp cơ sở dữ liệu hay shell; mọi hành
động đi qua tool schema, authorization và audit.

\subsection{Tool Calling}

Tool calling là cơ chế mô hình chọn hàm và tạo tham số có cấu trúc. Hệ thống phải
kiểm tra schema, quyền người dùng, phạm vi node và policy trước khi thực thi. Các
tool thay đổi trạng thái dùng hai pha: tạo kế hoạch để người dùng xem xét, sau đó
phê duyệt bằng định danh kế hoạch và token chống thực thi lại.

\subsection{MCP}

Model Context Protocol (MCP) chuẩn hóa cách AI application khám phá và gọi
tools/resources do server cung cấp. Trong phạm vi đồ án, MCP được dùng làm biên
tích hợp cho Knowledge và các công cụ vận hành. MCP không thay thế cơ chế xác
thực hoặc phân quyền của domain; server vẫn phải xác minh danh tính và chính
sách cho từng lần gọi.